\documentclass[a4paper]{article}
\usepackage{amsfonts,amsmath,amssymb,graphicx}

\begin{document}

\noindent \large Auteur: Michael Zielinski \\

\noindent \large Studierichting: Informatica\\

\noindent \large Oefeningenreeks 3E nr. 5 \\

\normalsize

\section*{Oplossing}

Stel het punt onder de Multanova als oorsprong en laat $x$ de afstand van de wagen tot die oorsprong zijn,\\
met $v=\dfrac{dx}{dt}$ de lineaire snelheid. \\

Omdat $\tan\theta = \dfrac{x}{20}$ geldt\\

$\theta = \operatorname{Bgtan}\left(\dfrac{x}{20}\right)$ \\

Differentieer naar de tijd:\\

$\omega = \dfrac{d\theta}{dt}
= \dfrac{d}{dx}\operatorname{Bgtan}\left(\dfrac{x}{20}\right)\dfrac{dx}{dt}
= \dfrac{20}{x^{2}+400}v$ \\

Hieruit volgt\\

$v = \dfrac{x^{2}+400}{20}\omega$ \\

Invullen van $x=15\ \mathrm{m}$ en $\omega=1,28\ \mathrm{s^{-1}}$ geeft\\

$$
v = \dfrac{15^{2}+400}{20}\cdot 1,28
= \dfrac{225+400}{20}\cdot 1,28
= 31,25\cdot 1,28
= 40\ \mathrm{m/s}
$$

Omzetten naar km/u:\\

$$
v = 40\cdot 3,6 = 144\ \mathrm{km/u}
$$

Dit is $144 - 120 = 24\ \mathrm{km/u}$ boven de toegelaten snelheid\\

en dus een zware overtreding. \\

\begin{thebibliography}{9}
% \bibitem{ref} Referentie
\end{thebibliography}

\end{document} 